\documentclass[11pt]{article}
\usepackage[a4paper,margin=25mm]{geometry}
\usepackage[UTF8,fontset=fandol,scheme=plain]{ctex}
\usepackage{amsmath,amssymb,amsthm,booktabs,array}
\usepackage[colorlinks=true,linkcolor=blue,urlcolor=blue]{hyperref}
\usepackage{xurl}
\newcommand{\Tr}{\operatorname{Tr}}
\newcommand{\Var}{\operatorname{Var}}
\newcommand{\Law}{\operatorname{Law}}
\newcommand{\E}{\mathbb E}
\newcommand{\lean}[1]{\nolinkurl{#1}}
\newcommand{\code}[1]{\texttt{\detokenize{#1}}}
\newtheorem{proposition}{Proposition}
\renewcommand{\refname}{References}
\setlength{\emergencystretch}{2em}
\title{The quadratic trace is deterministic:\\a disproof of Conjecture 00000007744}
\author{C0ldSmi1e}
\date{5 October 2026}
\begin{document}
\maketitle

\begin{abstract}
In the uniform simple $d$-regular graph model, the squared adjacency matrix has
trace $nd$ for every graph on $n$ vertices. Thus its centered quadratic trace,
with any deterministic matrix and trace normalization and with the stated
$\sqrt n$ factor, is identically zero. Its law is exactly $\delta_0$ at every
admissible size, and every weak limit has variance zero. For $d=3$, the
conjecture instead prescribes variance $8/9$. We formalize the actual uniform
model, polynomial traces, finite-dimensional laws, and the weak-limit
contradiction in Lean 4.19.0 with Mathlib v4.19.0.
\end{abstract}

\section{The original claim and the clause refuted}
The supplied English statement reads:
\begin{quote}\small
Definition: The spectral empirical measure of random d-regular graphs (uniform
model) converges to the Kesten-McKay law $KM_d$; the second-order limit is the
joint limit process of normalized trace polynomials times $\sqrt n$.
Conjecture: The limit is Gaussian with covariance $\Sigma(p,q)$ given by
$(d+1)$-regular constrained counts of noncrossing return path pairings of $p$
and $q$; for $p=q=x^2$ one has $\Var=4-12/d+8/d^2$, degenerating to the Wigner
value $4$ of GUE fluctuations as $d$ grows.
(Kesten-McKay second-order fluctuations)
\end{quote}
The Chinese version makes the centering explicit:
\begin{quote}\small
定义：随机 $d$-正则图（一致模型）的谱经验测度收敛到 Kesten--McKay 律
$KM_d$；其二阶极限分布指规范化迹多项式
$(\operatorname{tr}p(A)-\E)\cdot\sqrt n$ 的联合极限过程。
猜想：该极限为高斯族，协方差函数 $\Sigma(p,q)$ 由 $p,q$ 的非交叉返回路径配对的
$(d+1)$-正则约束计数给出；特别对 $p=q=x^2$，
$\Var=4-12/d+8/d^2$，且当 $d\to\infty$ 时退化为 GUE 波动的 Wigner 值 $4$。
\end{quote}
Both versions require the displayed quadratic-coordinate variance. We disprove
that explicit requirement at fixed degree $3$. We do not dispute the
first-order Kesten--McKay convergence or claim that every other coordinate
vanishes. No interpretation of the unspecified general pairing rule is needed.
The exact original Markdown file is included as \lean{original_conjecture.md}.

\section{The whole uniform model and its quadratic trace}
We use the conventional uniform model of finite labelled \emph{simple},
undirected $d$-regular graphs. Thus loops and parallel edges are absent; we do
not assert a result about an unconditioned multigraph configuration model.
The source states no connectedness requirement. On any finite label set $V$
of size $n$, let
\[
 \mathcal G_{V,d}=\{G:\ G\text{ is a simple graph on }V,
                   \ \deg_G(v)=d\text{ for every }v\in V\}.
\]
When this set is nonempty, give \emph{every} graph mass
$1/|\mathcal G_{V,d}|$. The Lean definitions use the subtype of
\code{SimpleGraph V} satisfying \code{IsRegularOfDegree d}, and
\lean{PMF.uniformOfFintype} on that entire subtype. An explicit theorem verifies
the equal masses. The measurable space consists of all subsets.

Write $A_G$ for the $0$--$1$ adjacency matrix. Let $a,b\in\mathbb R$ be
deterministic scalars and put
\[
 T_{a,b}(G)=b\Tr\bigl((aA_G)^2\bigr),\qquad
 Y_{a,b}(G)=\bigl(T_{a,b}(G)-\E T_{a,b}\bigr)\sqrt n.
\]
Ordinary adjacency and empirical normalized trace correspond to $a=1$ and
$b=1/n$. Allowing arbitrary $a,b$ also covers ordinary unnormalized trace and
deterministic rescalings of the adjacency matrix.

\begin{proposition}\label{prop:zero}
For every graph in $\mathcal G_{V,d}$,
\[
 T_{a,b}(G)=ba^2nd,\qquad \E T_{a,b}=ba^2nd,\qquad Y_{a,b}(G)=0.
\]
Consequently $\Law(Y_{a,b})=\delta_0$.
\end{proposition}
\begin{proof}
For each $v\in V$, symmetry and $0$--$1$ adjacency entries give
\[
 (A_G^2)_{vv}=\sum_{w\in V}(A_G)_{vw}(A_G)_{wv}
             =\sum_{w\in V}(A_G)_{vw}=\deg_G(v)=d.
\]
Summing the diagonal yields $\Tr(A_G^2)=nd$, and scalar multiplication gives
the asserted formula for $T_{a,b}$. It is a constant on the \emph{whole}
sample space, so its integral against the uniform probability measure is the
same constant. Subtracting that expectation gives zero pointwise. Pushing the
probability measure forward by this constant map gives $\delta_0$.
\end{proof}

\section{Nonempty large models and the actual weak limit}
For each $k\geq0$ take
\[
 V_{d,k}=\{0,\ldots,k\}\times\{0,\ldots,d\},\qquad
 n_{d,k}=(k+1)(d+1).
\]
The graph with $(i,u)$ adjacent to $(j,v)$ exactly when $i=j$ and $u\ne v$
is a disjoint union of $k+1$ copies of $K_{d+1}$, hence is $d$-regular.
This proves nonemptiness of the sample space at every such size. The random
graph is still drawn uniformly from \emph{all} regular graphs on $V_{d,k}$;
the block graph is only an existence witness. Lean proves the stated
cardinality, its positivity, and $n_{d,k}\to\infty$.

Let $a_k,b_k$ be any deterministic real sequences and let $\mu_k$ be the
actual law of $Y_{a_k,b_k}$ on this model. Proposition~\ref{prop:zero} gives
$\mu_k=\delta_0$ for every $k$, so $\mu_k\Rightarrow\delta_0$. If also
$\mu_k\Rightarrow\nu$, uniqueness of weak limits gives $\nu=\delta_0$, and
therefore
\[
 \Var_{\nu}(\operatorname{id})=0.
\]
This argument identifies the entire limiting law. It does not interchange
variance and weak limits. Lean uses the usual weak topology on
\lean{ProbabilityMeasure}\,$\mathbb R$, its Hausdorff property, and the
ordinary measure-theoretic variance.

For $d=3$ this is an unbounded sequence of admissible sizes $4(k+1)$. A limit
asserted over all admissible sizes must also hold along this sequence.
The finite trace identity holds on every finite label set; the product labels
impose no additional restriction on the sampled graphs.

\section{From the joint polynomial process to the contradiction}
For an arbitrary real polynomial $p$, the formal project defines
\[
 T_p(G)=b\Tr\bigl(p(aA_G)\bigr),\qquad
 Z_p(G)=\bigl(T_p(G)-\E T_p\bigr)\sqrt n
\]
using \lean{Polynomial.aeval}, the actual matrix trace and the actual uniform
integral. These observables are integrable on the finite sample space.
It proves that $p=x^2$ specializes to the preceding quadratic observable.
For any finite list $p_1,\ldots,p_m$, it constructs the genuine joint law of
$(Z_{p_1},\ldots,Z_{p_m})$ and identifies each coordinate marginal.
Continuous coordinate projection sends a weakly convergent sequence of these
joint laws to a weakly convergent sequence of their marginals.

The predicate \lean{RequiredSecondOrderLimit} states a \emph{necessary part}
of the original claim: these finite-dimensional laws have weak limits and
the one-coordinate list $p=x^2$ has the prescribed limiting variance.
Gaussianity and the other covariance prescriptions are further requirements
of the original claim. They are not silently redefined in this predicate.
The formal disproof of this necessary part is sufficient to refute the
original conjunction.

At $d=3$ the required variance is
\[
 4-\frac{12}{3}+\frac{8}{3^2}=\frac89\ne0.
\]
It is incompatible with the uniquely forced law $\delta_0$. Thus no proposed
Gaussian family with the source's covariance can be the joint limit.
A degenerate Gaussian coordinate would be consistent with zero fluctuations;
the contradiction is the explicitly positive variance.

The general formal conclusion is
\[
 \forall a,b:\mathbb N\to\mathbb R,\quad
 \neg\,\mathrm{RequiredSecondOrderLimit}(3,a,b),
\]
proved by \lean{no_required_second_order_limit}. The final theorem
\lean{conjecture7744_false} specializes to $a_k=1$ and $b_k=1/n_{3,k}$.
No unproved first-order spectral limit or central limit theorem is used.

\section{Formal and computational verification}
The project pins Lean 4.19.0 and Mathlib v4.19.0 (Mathlib commit
\lean{c44e0c8ee63ca166450922a373c7409c5d26b00b}), with all nine dependency
revisions recorded in \lean{lake-manifest.json}. The project build treats
warnings as errors. The author, coordinator and independent reviewer performed fresh builds.
The coordinator and reviewer also replayed every final module with
\code{-DwarningAsError=true}.

There are 49 explicitly authored declarations: 19 definitions, 26 theorems,
and four named instances. The full compiled inventory contains 83 declarations:
75 logical declarations and eight compiler-generated runtime wrappers.
Every logical declaration has only the standard Lean axioms
\lean{propext}, \lean{Classical.choice}, and \lean{Quot.sound}, or none.
There are no axiom declarations, admissions, \lean{native_decide}, partial
logical declarations, or unsafe/missing logical dependencies. The runtime
wrappers are retained in the inventory and are absent from logical dependency
closures. The author source has no unsafe proof mechanism.

The supplied exact-arithmetic Python check enumerates every labelled simple
cubic graph on four and six vertices, computes the full matrix square, and
then computes uniform expectations and centered variances:
\begin{center}
\begin{tabular}{rrrrr}
\toprule
$n$ & Graph count & $\Tr(A^2)/n$ & Centered variance & Claimed limit variance\\
\midrule
4 & 1 & 3 & 0 & $8/9$\\
6 & 70 & 3 & 0 & $8/9$\\
\bottomrule
\end{tabular}
\end{center}
The reviewer independently replayed that script and an alternative complete
edge-mask enumeration, with matching results. These finite calculations are
supporting checks only; the asymptotic disproof is the general Lean argument.

The accompanying README gives reproduction commands and the evidence map.
\code{Audit.lean} prints axiom dependencies for every authored declaration;
\code{Check.lean} also prints definitions and theorem/instance types.
The independent verifier inventories all compiled declarations and their
transitive dependencies. Reports, logs, source identities and auxiliary code
are included. Early failed author attempts retain their commands, exit codes
and diagnostics, but their overwritten source versions were not preserved;
contemporaneous source snapshots begin at final verification. Final-source
replays are fully bound to their hashes. Local validation is distinct from
maintainer acceptance.

\begin{thebibliography}{9}
\bibitem{source} The Last Math Competition, Conjecture 00000007744,
\href{https://github.com/The-Last-Math-Competition/The-Last-Math-Competition/blob/ba46ce0997cbf232fe95e5dced37ff4d32e1eb94/conjectures/00000007744.md}{exact source at the submission base}.
\bibitem{mathlib} Mathlib v4.19.0,
\lean{SimpleGraph.AdjMatrix}, \lean{Probability.Distributions.Uniform},
\lean{Probability.Variance}, and \lean{MeasureTheory.Measure.ProbabilityMeasure}.
Exact pinned-source definition extracts are retained in the coordinator evidence.
\end{thebibliography}
\end{document}
